\section{International and Regional Governance}
\label{sec:8.7}
No single country can bind AI development to altruistic and antifascist values on its own, and no two regions have tried to do it the same way.

Sections~\ref{sec:8.7.1} and \ref{sec:8.7.2} cover cooperation from the bottom up: how research institutions build the working relationships cross-border AI research depends on, and how the technical community writes standards for itself without waiting for governments. Section~\ref{sec:8.7.3} asks what that kind of collaboration can and cannot fix once the underlying problem is climate, disease, poverty, or conflict instead of AI ethics as such. Sections~\ref{sec:8.7.4} and \ref{sec:8.7.5} turn to jurisdictions: what a handful of regions have actually legislated, how little of it harmonizes, and what adapting a system to a local context does and does not license. Sections~\ref{sec:8.7.6} through \ref{sec:8.7.9} then take the whole landscape at the scale of states — compute and export controls, weapons and the vendor's place in the chain, treaties and declarations, and the digital divide — and section~\ref{sec:8.7.10} asks what a democratic dividend is actually worth.
